\documentclass[12pt,a4paper]{article}
\usepackage[utf8]{inputenc}
\usepackage[T1]{fontenc}
\usepackage[brazilian]{babel}
\usepackage{amsmath,amssymb}
\usepackage[margin=2.2cm]{geometry}
\usepackage{tikz}
\usepackage{pgfplots}
\pgfplotsset{compat=1.18}
\usetikzlibrary{arrows.meta,positioning}
\usepackage{tcolorbox}
\tcbuselibrary{skins,breakable}
\usepackage{enumitem}
\usepackage{xcolor}
\usepackage{booktabs}
\usepackage{hyperref}
\hypersetup{colorlinks=true,linkcolor=blue!50!black,urlcolor=blue!50!black}

\newtcolorbox{definicao}[1]{colback=blue!5,colframe=blue!55!black,fonttitle=\bfseries,coltitle=white,title={#1},breakable,boxrule=0.8pt,arc=2mm}
\newtcolorbox{exemplo}[1]{colback=green!5,colframe=green!35!black,fonttitle=\bfseries,coltitle=white,title={#1},breakable,boxrule=0.8pt,arc=2mm}
\newtcolorbox{atencao}[1]{colback=red!5,colframe=red!55!black,fonttitle=\bfseries,coltitle=white,title={#1},breakable,boxrule=0.8pt,arc=2mm}
\newtcolorbox{dica}[1]{colback=yellow!12,colframe=orange!70!black,fonttitle=\bfseries,coltitle=black,title={#1},breakable,boxrule=0.8pt,arc=2mm}

\title{\vspace{-1.5cm}\textbf{Tema 3 — Potenciação e Radiciação}\\[2pt]\large Regras que você usa em quase toda conta da prova}
\author{}
\date{}

\begin{document}
\maketitle
\thispagestyle{empty}

\section*{A ideia central}
Potência é \textbf{multiplicação repetida}: $a^n = \underbrace{a\cdot a\cdots a}_{n\text{ vezes}}$.
Raiz é a operação \textbf{inversa}: ``desfaz'' a potência.
\[
\sqrt[n]{a} = b \quad\Longleftrightarrow\quad b^n = a
\]
Tudo que vem a seguir são consequências lógicas dessas duas frases — não
decore, \textbf{deduza}.

\section{Potenciação}

\subsection{Expoentes inteiros}

\begin{definicao}{Expoente natural, zero e negativo}
\begin{align*}
a^n &= \underbrace{a \cdot a \cdots a}_{n} \quad (n \in \mathbb{N}^*)\\[4pt]
a^0 &= 1 \quad (a \neq 0)\\[4pt]
a^{-n} &= \frac{1}{a^n} \quad (a \neq 0)
\end{align*}
\end{definicao}

\begin{dica}{Por que $a^0=1$ e não $0$?}
Pense na sequência $a^3, a^2, a^1, a^0, a^{-1}$. Cada vez que o expoente cai
$1$, você \textbf{divide} por $a$:
\[
a^3 \xrightarrow{\div a} a^2 \xrightarrow{\div a} a^1 \xrightarrow{\div a} a^0 \xrightarrow{\div a} a^{-1}
\]
Com $a=2$: $8 \to 4 \to 2 \to 1 \to \tfrac12$. Para o padrão continuar, $a^0$
\textbf{tem} que ser $1$, e $a^{-1}$ \textbf{tem} que ser $\tfrac1a$. O
expoente negativo não é ``número negativo'', é apenas ``inverte a fração''.
\end{dica}

\begin{exemplo}{Aplicando direto}
\begin{align*}
3^3 &= 27\\
\left(\tfrac{3}{2}\right)^3 &= \tfrac{27}{8}\\
(-4)^4 &= 256 \quad\text{(expoente par: sinal negativo desaparece)}\\
\frac{1}{4^{-3}} &= 4^3 = 64 \quad\text{(inverter o inverso volta ao original)}\\
\frac{2^{-3}}{2^{-2}} &= 2^{-3-(-2)} = 2^{-1} = \tfrac12
\end{align*}
\end{exemplo}

\begin{atencao}{Cuidado com o sinal}
$(-4)^4 = 256$ (positivo), mas $-4^4 = -256$ (negativo)! O parêntese muda
tudo: sem parêntese, o sinal de menos \textbf{não} é elevado à potência.
Isso aparece na lista: $-2^2-(-5)^0 = -4 - 1 = -5$.
\end{atencao}

\subsection{As 5 propriedades operatórias}

\begin{center}
\renewcommand{\arraystretch}{1.6}
\begin{tabular}{@{}ll@{}}
\toprule
\textbf{Propriedade} & \textbf{Por que funciona (contando fatores)} \\
\midrule
$a^m \cdot a^n = a^{m+n}$ & $(a\cdots a)_{m} \cdot (a \cdots a)_n$ tem $m+n$ fatores \\
$\dfrac{a^m}{a^n} = a^{m-n}$ & cada $a$ de baixo cancela um de cima \\
$(a^m)^n = a^{m\cdot n}$ & $n$ grupos de $m$ fatores $= m\cdot n$ fatores \\
$(a\cdot b)^n = a^n \cdot b^n$ & distribui o expoente no produto \\
$\left(\dfrac{a}{b}\right)^n = \dfrac{a^n}{b^n}$ & distribui o expoente na fração \\
\bottomrule
\end{tabular}
\end{center}

\begin{exemplo}{Usando as propriedades (estilo da lista)}
\[
\left(\frac35\right)^5\cdot\left(\frac35\right)^{-6}\cdot\left(\frac35\right)^3
= \left(\frac35\right)^{5-6+3} = \left(\frac35\right)^{2} = \frac{9}{25}
\]
\[
\frac{7^3\cdot 7^4}{7^6} = \frac{7^{7}}{7^{6}} = 7^{1} = 7
\]
\[
2^3+\left(\frac12\right)^{-2} = 8 + 2^2 = 8+4=12
\]
\end{exemplo}

\begin{atencao}{O erro mais comum da prova}
$a^m \cdot a^n \neq a^{m\cdot n}$! Multiplicar potências de mesma base
\textbf{soma} os expoentes, não multiplica. Só se multiplica o expoente
quando é ``potência de potência'': $(a^m)^n = a^{mn}$.
\end{atencao}

\section{Radiciação}

\begin{definicao}{Raiz $n$-ésima}
\[
\sqrt[n]{a} = b \iff b^n = a, \qquad (\text{se } n \text{ par}, \text{ exige-se } a\ge 0 \text{ e } b\ge 0)
\]
Quando $n=2$ não escrevemos o índice: $\sqrt{a}=\sqrt[2]{a}$.
\end{definicao}

\begin{atencao}{Por que índice par exige radicando $\ge 0$?}
Em $\mathbb{R}$, todo número elevado a um expoente \textbf{par} é
$\ge 0$ (positivo $\times$ positivo, ou negativo $\times$ negativo, sempre
dá positivo). Logo não existe número real $b$ com $b^2 = -4$: não há raiz
quadrada real de número negativo. Já para índice \textbf{ímpar}
($\sqrt[3]{-8}=-2$, pois $(-2)^3=-8$) não há esse problema.
\end{atencao}

\subsection{Raiz como expoente fracionário — a ponte mais importante do tema}

\begin{definicao}{Convertendo raiz em potência}
\[
\sqrt[n]{a^m} = a^{\frac{m}{n}}
\]
Essa igualdade permite tratar \textbf{toda} raiz com as 5 propriedades de
potência que você já sabe — não precisa decorar regras novas de raiz.
\end{definicao}

\begin{exemplo}{Simplificando com expoente fracionário}
\[
\frac{x^{\frac13}\cdot x^{\frac14}}{x^{\frac34}} = x^{\frac13+\frac14-\frac34} = x^{\frac{4+3-9}{12}} = x^{-\frac{2}{12}} = x^{-\frac16} = \frac{1}{\sqrt[6]{x}}
\]
\[
\sqrt[3]{y^{12}} = y^{\frac{12}{3}} = y^{4}
\]
\[
\sqrt[3]{x^5 y^6} = (x^5y^6)^{\frac13} = x^{\frac53}y^{2} = x^{\frac53}y^2
\]
\end{exemplo}

\subsection{Simplificando radicais (fatoração dentro da raiz)}

\begin{exemplo}{$\sqrt{108x^8z^{18}}$}
Fatore o número e agrupe os expoentes em múltiplos do índice ($2$):
\[
108 = 4\cdot27 = 2^2\cdot 3^3
\]
\[
\sqrt{2^2\cdot3^3\cdot x^8 \cdot z^{18}} = \sqrt{2^2\cdot3^2\cdot3\cdot x^8\cdot z^{18}} = 2\cdot3\cdot x^4\cdot z^9\cdot\sqrt{3} = 6\sqrt3\,x^4z^9
\]
\textbf{Ideia}: tudo que tem expoente múltiplo do índice ``sai'' da raiz;
o que sobra fica dentro.
\end{exemplo}

\subsection{Racionalização — tirando raiz do denominador}

\begin{dica}{Por que racionalizar?}
Por convenção (e para facilitar comparação/soma de frações), não deixamos
raiz no denominador. A ideia sempre é: multiplique em cima e embaixo por
algo que \textbf{complete} a raiz do denominador para virar um número
inteiro, sem alterar o valor da fração (multiplicar por $\frac{k}{k}=1$).
\end{dica}

\begin{exemplo}{Caso 1 — denominador com raiz quadrada simples}
\[
\frac{6}{\sqrt2} = \frac{6}{\sqrt2}\cdot\frac{\sqrt2}{\sqrt2} = \frac{6\sqrt2}{2} = 3\sqrt2
\]
\end{exemplo}

\begin{exemplo}{Caso 2 — denominador com raiz de índice $n>2$}
Para $\sqrt[3]{10}$, falta ``completar'' o expoente de $10$ até $3$ (ele
tem expoente $1$, falta $2$):
\[
\frac{1}{\sqrt[3]{10}} = \frac{1}{10^{1/3}}\cdot\frac{10^{2/3}}{10^{2/3}} = \frac{10^{2/3}}{10} = \frac{\sqrt[3]{100}}{10}
\]
\end{exemplo}

\begin{exemplo}{Caso 3 — denominador com soma/diferença de raízes (usa produto notável)}
Multiplique pelo \textbf{conjugado} (mesma expressão, sinal trocado), pois
$(a-b)(a+b)=a^2-b^2$ elimina as raízes:
\[
\frac{2}{\sqrt5-\sqrt3}\cdot\frac{\sqrt5+\sqrt3}{\sqrt5+\sqrt3} = \frac{2(\sqrt5+\sqrt3)}{5-3} = \frac{2(\sqrt5+\sqrt3)}{2} = \sqrt5+\sqrt3
\]
\[
\frac{9}{4+\sqrt5}\cdot\frac{4-\sqrt5}{4-\sqrt5} = \frac{9(4-\sqrt5)}{16-5} = \frac{9(4-\sqrt5)}{11} = \frac{36-9\sqrt5}{11}
\]
\end{exemplo}

\section{Enxergando raiz e potência como funções inversas}

O gráfico deixa claro por que ``raiz desfaz potência'': as curvas
$y=x^2$ (para $x\ge0$) e $y=\sqrt{x}$ são \textbf{reflexos} uma da outra em
torno da reta $y=x$.

\begin{center}
\begin{tikzpicture}
\begin{axis}[
  width=9cm, height=9cm,
  xmin=-0.3, xmax=4.5, ymin=-0.3, ymax=4.5,
  axis lines=middle,
  xlabel=$x$, ylabel=$y$,
  samples=100, thick,
  legend pos=north west,
  legend style={font=\small}
]
\addplot[domain=0:2.12,blue] {x^2};
\addlegendentry{$y=x^2\ (x\ge0)$}
\addplot[domain=0:4.5,red] {sqrt(x)};
\addlegendentry{$y=\sqrt{x}$}
\addplot[domain=0:4.3,dashed,gray] {x};
\addlegendentry{$y=x$ (espelho)}
\end{axis}
\end{tikzpicture}
\end{center}

\section*{Resumo relâmpago}
\begin{itemize}[nosep]
  \item $a^{-n} = \dfrac{1}{a^n}$; \quad $a^0=1$.
  \item Multiplicar mesma base: \textbf{soma} expoente. Potência de potência: \textbf{multiplica} expoente.
  \item $\sqrt[n]{a^m} = a^{m/n}$ — converta raiz em potência sempre que travar.
  \item Índice par $\Rightarrow$ radicando $\ge 0$.
  \item Racionalizar: multiplique por $1$ disfarçado (a própria raiz, o que falta para o índice, ou o conjugado).
\end{itemize}

\end{document}
